% =====================================================================
%  Conferenze — Evento 2: On The Road, Zagarolo, giorno 2
%  Introduzione alla giornata e mappa dei capitoli C3.x.
%  Capitolo non numerato. Nessun pacchetto aggiuntivo.
% =====================================================================

\chapter*{On The Road, Zagarolo: secondo giorno}
\phantomsection
\addcontentsline{toc}{chapter}{On The Road, Zagarolo: secondo giorno}

Il secondo giorno di \textbf{On The Road} riprende molti fili del primo: i
relatori si citano a vicenda (Jonas di Raw Fury, Manos, Tony), e i temi si
spostano dall'ingresso nell'industria al \textbf{come si lavora} una volta dentro:
gli eventi a cui partecipare, il debito tecnico, il lavoro invisibile del producer,
la scelta di fondare uno studio, e infine il rapporto tra videogiochi e giochi da
tavolo.

Gli interventi registrati sono cinque:
\begin{enumerate}
  \item l'\textbf{importanza degli eventi internazionali} per gli studi
        indipendenti, con il London Games Festival
        (capitolo~\ref{cap:c3-1-eventi-internazionali});
  \item il \textbf{debito tecnico}: riconoscerlo, prevenirlo, gestirlo
        (capitolo~\ref{cap:c3-2-debito-tecnico});
  \item \textbf{People Over Processes}, il lavoro invisibile del producer
        (capitolo~\ref{cap:c3-3-people-over-processes});
  \item \textbf{dal tripla A all'indie}: due veterani raccontano la fondazione del
        loro studio (capitolo~\ref{cap:c3-4-day-4-night});
  \item \textbf{giochi da tavolo e videogiochi}: testare un gioco di carte con il
        machine learning e adattare un videogioco al tavolo
        (capitolo~\ref{cap:c3-5-giochi-da-tavolo}).
\end{enumerate}
